\section{About the Topic}

\subsection{Key Terms}

\textbf{Illicit financial flows (IFFs).} The agenda’s headline concept. The first UN-wide definition, set jointly by UNCTAD and UNODC, describes them as financial flows that are illicit in origin, transfer, or use and that cross country borders, and they are a target of Sustainable Development Goal 16.4.\footnote{United Nations Conference on Trade and Development \& United Nations Office on Drugs and Crime. (2020). \emph{Conceptual framework for the statistical measurement of illicit financial flows}. United Nations. \url{https://unctad.org/publication/conceptual-framework-statistical-measurement-illicit-financial-flows}}

\textbf{Money laundering.} The process of disguising the criminal origin of money so it can be used freely. The UN’s 1988 Vienna Convention defines it as converting or transferring property to conceal its illicit source, and it classically moves through three stages: placement (introducing the proceeds into the financial system), layering (moving them through transactions to break the trail), and integration (returning them to the owner as apparently legitimate wealth).\footnote{United Nations Office on Drugs and Crime. (n.d.). \emph{Money laundering: Overview}. \url{https://www.unodc.org/unodc/en/money-laundering/overview.html}}

\textbf{Predicate offence.} The underlying crime that generates the money in the first place, such as drug or arms trafficking, fraud, or corruption. Laundering is the secondary act of cleaning those proceeds.\footnote{Financial Action Task Force. (2023). \emph{International standards on combating money laundering and the financing of terrorism and proliferation: The FATF recommendations}. FATF. \url{https://www.fatf-gafi.org/en/publications/Fatfrecommendations/Fatf-recommendations.html}}

\textbf{Tax haven and secrecy jurisdiction.} Related but distinct. A tax haven offers very low or zero taxation, often used to shift corporate profits; a secrecy jurisdiction offers legal and financial secrecy to people and entities based elsewhere, concealing who owns what. The Tax Justice Network treats secrecy as a spectrum on which every jurisdiction sits to some degree, rather than a fixed blacklist.\footnote{Tax Justice Network. (n.d.-a). \emph{Tax havens and secrecy jurisdictions}. \url{https://taxjustice.net/topics/tax-havens-and-secrecy-jurisdictions/}}

\textbf{Offshore financial centre.} A jurisdiction that provides financial services to non-residents on a scale far larger than its own economy, typically pairing light regulation with strong confidentiality.

\textbf{Shell company.} A legal entity with no real operations or employees, used to hold assets or move funds while obscuring the people behind it; shells can be stacked in chains across jurisdictions.

\begin{itemize}
  \item \textbf{Beneficial owner.} The real, flesh-and-blood person who ultimately owns or controls an entity, as opposed to the nominee or company named on paper. Concealing beneficial ownership is the core technique of offshore secrecy.\footnote{FATF (2023), \emph{The FATF recommendations}.}
  \item \textbf{Asset recovery.} The process of identifying, freezing, confiscating, and returning the proceeds of crime and corruption to their rightful owners, addressed in Chapter V of the UN Convention against Corruption.\footnote{United Nations. (2004). \emph{United Nations Convention against Corruption}, ch. V. \url{https://www.unodc.org/unodc/en/corruption/uncac.html}}
  \item \textbf{Financial sovereignty.} The authority of a state to govern its own monetary, fiscal, and financial affairs, including its laws on banking, taxation, corporate registration, and the movement of capital across its borders. It is the principle on which both legitimate financial privacy and harmful secrecy rest, and the friction it creates with cross-border enforcement is the central problem of this committee.
\end{itemize}

\subsection{Introduction to the Topic}

The defining dimension of this topic that we would want to explore in the committee is the tension between sovereignty and accountability. Under international law, each state decides for itself how to tax, how tightly to regulate its banks, and how much of a company’s ownership to make public, and when a jurisdiction chooses minimal taxation and strong secrecy it is exercising that sovereign right. The difficulty is that financial secrecy is rarely consumed at home; it is sold to non-residents.\footnote{Tax Justice Network. (n.d.-b). \emph{About the Financial Secrecy Index}. \url{https://fsi.taxjustice.net/about-the-index/}} A secrecy law passed in one country becomes, in practice, an escape route for the citizens, companies, and criminals of every other country, letting them place wealth beyond the reach of the authorities that would otherwise tax or investigate it. One state’s sovereignty thus becomes another state’s enforcement gap, and there is no global authority that can simply close it.

The stakes are large and concrete. UNODC estimates that between two and five per cent of global GDP, somewhere between 800 billion and 2 trillion US dollars, is laundered every year, and the same secrecy that hides a tax evader’s fortune also conceals the proceeds of corruption, fraud, and the trafficking of drugs and arms.\footnote{UNODC (n.d.), \emph{Money laundering: Overview}.} Money stolen from public budgets is moved offshore and lost to the populations it was meant to serve, and criminal organizations rely on opaque companies and accounts to move and store their earnings. Because these flows cross borders while investigators remain confined within them, no single state can follow the money to its end. That structural mismatch, borderless money set against border-bound enforcement, is what makes international cooperation indispensable and also what makes it so difficult to achieve without touching the very sovereignty that members guard.

For delegates, then, the task is not to abolish financial privacy or to override national law, neither of which is realistic or desirable, but to find mechanisms that let states cooperate against the criminal abuse of the financial system while respecting their differing legal traditions and their legitimate interests in confidentiality. Holding that balance, between the accountability the international community demands and the sovereignty states will not surrender, is the work of this committee.

\subsection{History of the Topic}

The architecture of financial secrecy is older than the crimes it now conceals. Swiss bankers had cultivated client confidentiality since the eighteenth century, and Switzerland codified it in the Federal Act on Banks and Savings Banks of 1934, which made disclosing a client’s identity a federal crime. Political instability, war, and rising taxes across early-twentieth-century Europe drove a flood of capital into Swiss accounts, and the combination of secrecy, a stable currency, and neutrality turned the country into the world’s preeminent haven for hidden wealth.\footnote{Palan, R., Murphy, R., \& Chavagneux, C. (2010). \emph{Tax havens: How globalization really works}. Cornell University Press.}

After the Second World War, the model spread. As trade and capital globalized, a network of offshore financial centres grew up across the Caribbean, the European microstates, and a number of British and Crown dependencies, each offering some mix of low tax, light regulation, and confidentiality to non-residents.\footnote{Palan, Murphy \& Chavagneux (2010), \emph{Tax havens}.} The essential instrument was the shell company, an entity that exists only on paper and can be layered in chains across several jurisdictions until the real owner disappears from view, and by the late twentieth century an entire professional industry of lawyers, accountants, and corporate service providers had grown up to build and maintain these structures.

International concern hardened as the link between secrecy and serious crime became undeniable. The UN’s 1988 Vienna Convention criminalized the laundering of drug proceeds, the G7 created the Financial Action Task Force in 1989 to set anti-money-laundering standards, and in 2015 the international community wrote the reduction of illicit financial flows into the Sustainable Development Goals.\footnote{FATF (2023), \emph{The FATF recommendations}; UNODC (n.d.), \emph{Money laundering: Overview}; UNCTAD \& UNODC (2020), \emph{Conceptual framework}.} These milestones marked a shift from treating secrecy as a private matter to treating it as a shared problem of security and development, though the detailed record of what has since been done about it belongs to the sections that follow.

What turned the issue from a technical debate into a public reckoning was a series of journalistic leaks. Beginning with the International Consortium of Investigative Journalists’ Offshore Leaks in 2013 and culminating in the Panama Papers of 2016 and the Pandora Papers of 2021, millions of leaked records laid bare the inner workings of the offshore world. The Panama Papers alone comprised some 11.5 million documents from a single law firm and exposed the secret holdings of twelve world leaders along with scores of other politicians and criminals;\footnote{International Consortium of Investigative Journalists. (2016). \emph{The Panama Papers: Exposing the rogue offshore finance industry}. \url{https://www.icij.org/investigations/panama-papers/}} the Pandora Papers drew on nearly twelve million records from fourteen providers and named more than 330 politicians and 35 current or former heads of state across more than ninety countries.\footnote{International Consortium of Investigative Journalists. (2021). \emph{Pandora Papers}. \url{https://www.icij.org/investigations/pandora-papers/}} The leaks showed that the same machinery served tax-shy billionaires, corrupt officials, and organized crime alike, and that the most significant secrecy enablers were not only small islands but included some of the world’s largest economies.\footnote{Tax Justice Network (n.d.-b), \emph{About the Financial Secrecy Index}.}

The exposure accelerated a retreat from absolute secrecy. Under pressure from the OECD and the G20, dozens of jurisdictions adopted the automatic exchange of financial account information, and Switzerland, the historic home of bank secrecy, began exchanging client data with foreign tax authorities in 2018.\footnote{Organisation for Economic Co-operation and Development. (2014). \emph{Standard for automatic exchange of financial account information in tax matters}. OECD Publishing. \url{https://www.oecd.org/tax/automatic-exchange/}} Yet the system has proved adaptable rather than defeated: as some doors closed, business migrated toward jurisdictions that stayed opaque, including, by several measures, certain states within the United States itself.\footnote{Tax Justice Network (n.d.-b), \emph{About the Financial Secrecy Index}.} The result is the landscape this committee inherits, one in which a formal norm of transparency now coexists with a resilient infrastructure of secrecy, and in which the gap between what states promise and what they enforce remains wide.
